2. \underline{Numerical results}

2.1. Let $\underline a=(a_i)_{1\leq i\leq d}$ be a sequence of integers
$\geq1$, and let $V_n(\underline a)$ denote a smooth complete intersection
of dimension $n$ and multidegree $\underline a$ in a projective space
$\mathbb P^{n+d}$ over a field $k$.  Put

\begin{equation*}
h^{p,q}(\underline a)
=\dim_k H^q\bigl(V_{p+q}(\underline a),\Omega^p\bigr)
\end{equation*}

and

\begin{equation*}
h_0^{p,q}(\underline a)=h^{p,q}(\underline a)-\delta_{p,q}.
\end{equation*}

Thus the $h_0^{p,q}(\underline a)$ are the Hodge numbers of the
middle-dimensional primitive cohomology of a complete intersection of
multidegree $\underline a$.  They depend only on $\underline a$, $p$, and
$q$, by 1.5.  Set

\begin{equation}\tag{2.1.1}
H(\underline a)=\sum_{p,q\geq0}h_0^{p,q}(\underline a)y^pz^q
\in\mathbb Z[[y,z]].
\end{equation}

For $d=1$ and $\underline a=(a)$, put
$h^{p,q}(a)=h^{p,q}(\underline a)$ and $H(a)=H(\underline a)$.  For $d=0$,
put $\sup(a_i)=1$.

The following formula is the special case $k=0$ of formula (2), p.~160,
in Hirzebruch [3]:

%%%%%%%%%% source scan idx 59 | folio 52 %%%%%%%%%%

\begin{equation}\tag{2.2}
\sum_{n=0}^{\infty}\sum_p
\chi\bigl(V_n(\underline a),\Omega^p\bigr)y^pz^{n+d}
=\frac{1}{(1+yz)(1-z)}
\prod_{i=1}^d
\frac{(1+yz)^{a_i}-(1-z)^{a_i}}
{(1+yz)^{a_i}z+y(1-z)^{a_i}}.
\end{equation}

We shall transform this into the following formula.

\underline{Theorem} 2.3 (Hirzebruch). \underline{With the notation of 2.1,
one has}

\begin{equation*}
H(\underline a)=\frac{1}{(1+y)(1+z)}\left[
\prod_{i=1}^d
\frac{(1+y)^{a_i}-(1+z)^{a_i}}
{-(1+y)^{a_i}z+(1+z)^{a_i}y}-1\right].
\end{equation*}

By 1.5,

\begin{equation*}
\chi\bigl(V_n(\underline a),\Omega^p\bigr)
=(-1)^p+(-1)^{n-p}h_0^{p,n-p}(\underline a).
\end{equation*}

Consequently, the left-hand side $I$ of (2.2) can be rewritten as

\begin{align*}
I
&=\sum_{n=0}^{\infty}\sum_{p+q=n}
\left((-1)^p+(-1)^q h_0^{p,q}(\underline a)\right)y^pz^{p+q+d}\\
&=\sum_{p,q\geq0}(-1)^p(yz)^pz^qz^d
+\sum_{p,q\geq0}(-1)^q h_0^{p,q}(\underline a)(yz)^pz^qz^d\\
&=z^d\left[\frac{1}{(1+yz)(1-z)}
+H(\underline a)(yz,-z)\right].
\end{align*}

The right-hand side $II$ of (2.2) can be rewritten as

\begin{equation*}
II=\frac{z^d}{(1+yz)(1-z)}
\prod_{i=1}^d
\frac{(1+yz)^{a_i}-(1-z)^{a_i}}
{(1+yz)^{a_i}z+yz(1-z)^{a_i}}.
\end{equation*}

After cancelling $z^d$ in $I=II$, one obtains

\begin{equation*}
H(\underline a)(yz,-z)
=\frac{1}{(1+yz)(1-z)}\left[
\prod_{i=1}^d
\frac{(1+yz)^{a_i}-(1-z)^{a_i}}
{(1+yz)^{a_i}z+(1-z)^{a_i}yz}-1\right],
\end{equation*}

and 2.3 follows by a change of variables.

%%%%%%%%%% source scan idx 60 | folio 53 %%%%%%%%%%

\underline{Corollary} 2.4. \underline{With the notation of 2.1:}

\begin{enumerate}[label=(\roman*),leftmargin=2.6em]
\item \underline{For $d=1$, one has}
\begin{equation*}
\begin{aligned}
H(a)
&=-\frac{(1+y)^{a-1}-(1+z)^{a-1}}
{(1+y)^az-(1+z)^ay}\\[1ex]
&=\frac{\displaystyle\sum_{i,j\geq0}
\binom{a-1}{i+j+1}y^iz^j}
{\displaystyle1-\sum_{i,j\geq1}\binom{a}{i+j}y^iz^j}.
\end{aligned}
\end{equation*}

\item \underline{Writing $|P|$ for the number of elements of a subset
$P$ of $[1,d]$, one has}
\begin{equation*}
H(\underline a)=
\sum_{\substack{P\subset[1,d]\\P\ne\varnothing}}
(1+y)^{|P|-1}(1+z)^{|P|-1}
\prod_{i\in P}H(a_i).
\end{equation*}
\end{enumerate}

By 2.3,

\begin{align*}
H(a)
&=-\frac{1}{(1+y)(1+z)}\left[
\frac{(1+y)^a-(1+z)^a}{(1+y)^az-(1+z)^ay}+1\right]\\
&=-\frac{1}{(1+y)(1+z)}\left[
\frac{(1+y)^a(1+z)-(1+z)^a(1+y)}
{(1+y)^az-(1+z)^ay}\right]\\
&=-\frac{(1+y)^{a-1}-(1+z)^{a-1}}
{(1+y)^az-(1+z)^ay}.
\end{align*}

Let $N$ and $D$ be the numerator and denominator in this last expression,
so that $H(a)=-N/D$.

%%%%%%%%%% source scan idx 61 | folio 54 %%%%%%%%%%

One has

\begin{equation*}
N=\sum_k\binom{a-1}{k}(y^k-z^k)
=(y-z)\sum_{i,j\geq0}\binom{a-1}{i+j+1}y^iz^j
\end{equation*}

and

\begin{align*}
D
&=\sum_k\binom ak(y^kz-yz^k)
=(z-y)+yz\sum_{k\geq1}\binom ak(y^{k-1}-z^{k-1})\\
&=(y-z)\left(-1+yz\sum_{i,j\geq0}
\binom a{i+j+2}y^iz^j\right)\\
&=(y-z)\left(-1+\sum_{i,j\geq1}
\binom a{i+j}y^iz^j\right).
\end{align*}

This proves (i).

Now invert the formula in 2.3 when $d=1$.  One finds

\begin{equation*}
-\frac{(1+y)^a-(1+z)^a}{(1+y)^az-(1+z)^ay}
=1+(1+y)(1+z)H(a).
\end{equation*}

Substituting this identity into 2.3 gives

\begin{align*}
H(\underline a)
&=\frac{1}{(1+y)(1+z)}\left[
\prod_{i=1}^d\left(1+(1+y)(1+z)H(a_i)\right)-1\right]\\
&=\frac{1}{(1+y)(1+z)}\left[
\sum_{\substack{P\subset[1,d]\\P\ne\varnothing}}
(1+y)^{|P|}(1+z)^{|P|}\prod_{i\in P}H(a_i)\right],
\end{align*}

which proves (ii).
